\chapter{Common Distributions}

\section{Uniform Distribution}
The simplest distribution is the \emph{discrete uniform} distribution.  It
describes a random variable that takes one of a finite set of values, all with
equal probability.  Examples:
\begin{itemize}
  \item a fair six-sided dice: each of \(1,2,3,4,5,6\) has probability \(1/6\);
  \item a single PowerSpin wheel result (ignoring the symbol grouping) when each
        of the 27 positions is equally likely.
\end{itemize}

If a random variable \(X\) is uniformly distributed over
\(\{1,2,\dots,n\}\), we write
\[
  P(X = k) = \frac{1}{n}, \quad k = 1,2,\dots,n.
\]
Its mean and variance are
\[
  \mathbb{E}[X] = \frac{n+1}{2}, \qquad
  \mathrm{Var}(X) = \frac{n^2-1}{12}.
\]

Uniform distributions capture the ideal of ``no preference'' among outcomes.
They are the starting point for most lottery and dice models in this book.

\section{Geometric Distribution}
The \emph{geometric} distribution describes the number of trials needed to get
the first \emph{success} in a sequence of independent trials, each with the
same success probability \(p\).  Typical stories:
\begin{itemize}
  \item ``How many tickets do I need to buy until I win once?''
  \item ``How many times do I roll until I see my favourite number?''
\end{itemize}

Let \(X\) be the number of trials needed to get the first success.  Then
\[
  P(X = k) = (1-p)^{k-1}p, \quad k = 1,2,3,\dots
\]
That is, we fail \(k-1\) times (each with probability \(1-p\)) and then
succeed on the \(k\)-th trial.

The mean and variance of a geometric random variable are
\[
  \mathbb{E}[X] = \frac{1}{p}, \qquad
  \mathrm{Var}(X) = \frac{1-p}{p^2}.
\]

The geometric distribution appears repeatedly in this book.  In the coupon
collector's problem, each \(t_i\) (the waiting time for the next new coupon) is
geometrically distributed with a suitable success probability \(p_i\).  In
lottery testing, geometric waiting times can describe how long we expect to
wait until a rare event occurs under a fair mechanism.

